\documentclass[
    12pt,
    a4paper,
    sumario=tradicional,
    chapter=TITLE
]{abntex2}

\usepackage[alf]{abntex2cite}
\usepackage[utf8]{inputenc}
\usepackage[T1]{fontenc}
\usepackage{graphicx}
\usepackage{booktabs}
\usepackage{longtable}
\usepackage{listings}
\usepackage{xcolor}
\usepackage{hyperref}
\usepackage{indentfirst}
\usepackage{microtype}

\hypersetup{
    colorlinks=true,
    linkcolor=black,
    citecolor=black,
    urlcolor=blue
}

\lstset{
    basicstyle=\ttfamily\small,
    keywordstyle=\color{blue},
    commentstyle=\color{gray},
    breaklines=true,
    frame=single,
    tabsize=2
}

\renewcommand{\ABNTEXchapterfont}{\fontfamily{ptm}\selectfont\bfseries}
\renewcommand{\ABNTEXsectionfont}{\fontfamily{ptm}\selectfont\bfseries}
\renewcommand{\ABNTEXsubsectionfont}{\fontfamily{ptm}\selectfont\bfseries}

\titulo{Cardly: aplicativo móvel de flashcards com revisão espaçada e integração cliente-servidor}
\autor{Hugo de Freitas Evangelista \and Paulo Henrique Calisto dos Santos \and Hudson Junior Rodrigues Silva \and Ivan Aloizio de Lana Filho}
\local{Governador Valadares -- MG}
\data{2026}
\instituicao{Centro Universitário do Leste de Minas Gerais -- UNILESTE}
\tipotrabalho{Trabalho de Conclusão de Disciplina (Projeto Final)}
\orientador{Prof. [Nome do orientador]}
\preambulo{Relatório técnico do projeto Cardly, desenvolvido na disciplina de Sistemas Móveis, apresentando arquitetura, requisitos, implementação e validação de um aplicativo de estudo com flashcards integrado a API REST em Spring Boot.}

\begin{document}

\frenchspacing
\imprimircapa
\imprimirfolhaderosto

\begin{resumo}
O Cardly é um aplicativo móvel multiplataforma inspirado no Anki, voltado ao estudo de diferentes assuntos por meio de flashcards e revisão espaçada. O front-end foi construído com React Native, Expo e TypeScript; o back-end, com Spring Boot e PostgreSQL. Este relatório documenta requisitos funcionais e técnicos, arquitetura cliente-servidor, camadas de segurança com JWT e Google SSO, estratégias de teste e preparação para deploy em produção. O texto está alinhado aos critérios de avaliação da disciplina e às normas ABNT para trabalhos acadêmicos.

\textbf{Palavras-chave}: flashcards; revisão espaçada; React Native; Spring Boot; arquitetura cliente-servidor; JWT.
\end{resumo}

\tableofcontents
\cleardoublepage

\chapter{Introdução}

\section{Contexto e motivação}

O aprendizado autodidata exige ferramentas que combinem organização de conteúdo, repetição inteligente e acompanhamento de progresso. Aplicativos de flashcards, como o Anki \cite{anki_manual}, consolidaram-se como referência nesse cenário ao aplicar técnicas de repetição espaçada, nas quais o intervalo entre revisões depende do desempenho do estudante \cite{anki_sm2,medium_spaced_repetition}.

O \textbf{Cardly} nasce nesse contexto como projeto acadêmico da disciplina de Sistemas Móveis, desenvolvido por estudantes do sétimo semestre de Engenharia de Software. Trata-se de um clone funcional inspirado no Anki, com foco em estudo mobile, dashboard de progresso, comunidade de assuntos públicos e painel administrativo. O sistema foi concebido para demonstrar integração completa entre front-end móvel (React Native com Expo e TypeScript) e back-end REST (Spring Boot com Java), conforme exigido pelo professor \cite{react_native_docs,expo_docs,spring_boot_docs}.

A proposta inicial foi apresentada ao docente por meio do documento \textit{Flashcards.pdf}, que orientou escopo, telas principais e regras de negócio. Este relatório expande essa proposta em documentação técnica formal, alinhada às normas ABNT \cite{abnt_nbr14724,abnt_nbr6023} e ao manual institucional UNILESTE utilizado pela equipe.

\section{Problema}

Estudantes universitários frequentemente acumulam material de diversas disciplinas sem um fluxo estruturado de revisão. Planilhas e anotações estáticas não registram dificuldade por carta, não impedem revisões prematuras e não oferecem visão temporal do hábito de estudo. Além disso, compartilhar baralhos entre colegas exige duplicação manual de conteúdo.

O problema abordado consiste em disponibilizar um aplicativo móvel que:
\begin{itemize}
    \item permita criar assuntos e cartas personalizadas;
    \item aplique revisão espaçada com intervalos definidos pelo sistema;
    \item exiba progresso e histórico de estudo em dashboard;
    \item ofereça assuntos públicos clonáveis pela comunidade;
    \item disponibilize gestão centralizada por perfil administrador.
\end{itemize}

\section{Solução proposta}

O Cardly resolve o problema por meio de uma arquitetura cliente-servidor \cite{tanenbaum2015,rest_fielding}: o aplicativo mobile consome uma API stateless protegida por JWT, persistindo dados em PostgreSQL \cite{postgresql_docs}. A lógica de dificuldade e agendamento de cartas permanece no back-end, garantindo consistência independentemente do cliente \cite{clean_architecture_martin}.

O front-end utiliza React Navigation para fluxos autenticados e não autenticados \cite{react_navigation}, NativeWind para estilização \cite{nativewind_docs,tailwind_docs} e animação de virada de carta acionada por toque, inspirada em referências de flip card \cite{w3schools_flip,reanimated_docs}. O back-end organiza-se em camadas Controller, Service, Repository e Specification \cite{spring_boot_docs,jpa_specification}, com migrations versionadas via Flyway \cite{flyway_docs}.

\section{Escopo acadêmico e produção}

Embora o projeto tenha caráter acadêmico e prazo limitado de operação em produção, o professor exige preparação para deploy real \cite{twelve_factor,vercel_docs,railway_docs}. A equipe adota variáveis de ambiente, build reprodutível e separação entre configuração de desenvolvimento e produção. Repositórios públicos concentram o código-fonte \cite{cardly_backend_repo,cardly_frontend_repo}.

\section{Organização deste relatório}

Este documento está estruturado da seguinte forma: o Capítulo~\ref{cap:objetivos} apresenta objetivos e mapeamento explícito aos critérios de avaliação; o Capítulo~\ref{cap:requisitos} detalha requisitos funcionais e técnicos; o Capítulo~\ref{cap:arquitetura} descreve a arquitetura integrada; os Capítulos~\ref{cap:backend} e~\ref{cap:frontend} aprofundam implementação; segurança, testes e deploy são tratados nos Capítulos~\ref{cap:seguranca},~\ref{cap:testes} e~\ref{cap:deploy}; por fim, o Capítulo~\ref{cap:conclusao} consolida resultados e trabalhos futuros.

\section{Justificativa do ecossistema mobile}

Diferentemente de uma aplicação web genérica, o Cardly explora estudo fragmentado no cotidiano do estudante: sessões curtas, feedback tátil imediato na virada de cartas e notificações futuras de revisão \cite{mobile_first_nielsen}. A experiência foi pensada para telas de smartphone, com gestos e componentes nativos via Expo \cite{expo_docs,youtube_expo_typescript}, justificando o uso do stack mobile exigido pela disciplina.

\chapter{Objetivos}
\label{cap:objetivos}

\section{Objetivo geral}

Desenvolver e documentar o Cardly, aplicativo móvel de flashcards com revisão espaçada, demonstrando integração completa entre React Native (Expo/TypeScript) e Spring Boot (Java), incluindo autenticação, CRUD, comunicação REST, organização em camadas, experiência de usuário responsiva e preparação para ambiente de produção.

\section{Objetivos específicos}

\begin{enumerate}
    \item Implementar fluxos de autenticação com credenciais locais e Google SSO, utilizando JWT stateless \cite{jwt_io,google_oauth,spring_security_jwt}.
    \item Modelar assuntos (\textit{decks}), cartas e respostas com regras de dificuldade e intervalos entre tentativas definidos pelo sistema.
    \item Disponibilizar dashboard com indicadores de progresso e calendário de dias com atividade de estudo, inspirado em aplicativos de acompanhamento diário \cite{clue_app}.
    \item Expor API REST paginada com endpoints \texttt{/search} baseados em Specification, sem \texttt{getAll} genérico \cite{jpa_specification}.
    \item Construir front-end com NativeWind, modais de confirmação e toasts equilibrados \cite{nativewind_docs,toast_message_rn}.
    \item Oferecer painel exclusivo de superadministrador para gestão de usuários, assuntos e cartas.
    \item Validar o sistema com testes automatizados e evidências visuais para arguição técnica \cite{junit5_spring,jest_testing_library}.
    \item Preparar pipeline de deploy (Railway/Vercel ou equivalente) para demonstração funcional \cite{railway_docs,vercel_docs,expo_eas_build}.
\end{enumerate}

\section{Mapeamento aos critérios de avaliação do professor}

A Tabela~\ref{tab:criterios-avaliacao} relaciona explicitamente os critérios divulgados na plataforma da disciplina (conforme \texttt{todo.md} do projeto) às seções deste relatório e às evidências implementadas no Cardly. Esse mapeamento orienta a apresentação prática e a arguição técnica.

\begin{table}[htbp]
    \caption{Mapeamento entre critérios de avaliação, seções do relatório e evidências no Cardly}
    \label{tab:criterios-avaliacao}
    \centering
    \small
    \begin{tabular}{p{4.2cm} p{3.8cm} p{6.5cm}}
        \toprule
        \textbf{Critério (disciplina)} & \textbf{Seção do relatório} & \textbf{Evidência no projeto} \\
        \midrule
        Integração Front-End (RN/Expo/TS) e Back-End (Spring Boot) &
        Cap.~\ref{cap:arquitetura}, \ref{cap:backend}, \ref{cap:frontend} &
        Módulos \texttt{*Api.ts} no app; controllers REST no backend; JSON paginado \\
        \addlinespace
        Arquitetura cliente-servidor &
        Cap.~\ref{cap:arquitetura} &
        Diagramas de camadas; API stateless; PostgreSQL centralizado \\
        \addlinespace
        Autenticação, CRUD e comunicação com API &
        Cap.~\ref{cap:seguranca}, \ref{cap:backend}, \ref{cap:frontend} &
        JWT + Google SSO; endpoints protegidos; CRUD de decks/cartas/usuários \\
        \addlinespace
        Organização de código e separação de responsabilidades &
        Cap.~\ref{cap:backend}, \ref{cap:frontend} &
        Padrão 1.0/1.1 no backend; screens, components, api e auth no front \\
        \addlinespace
        UI/UX, responsividade e experiência &
        Cap.~\ref{cap:frontend} &
        FlipCard animado; paleta 2.0; toasts e modais; dashboard \\
        \addlinespace
        Preparo para produção &
        Cap.~\ref{cap:deploy} &
        Variáveis de ambiente; build EAS; deploy Railway; CORS \\
        \addlinespace
        Funcionamento completo da aplicação &
        Cap.~\ref{cap:testes}, \ref{cap:conclusao} &
        Testes de serviço; roteiro de demonstração; prints \\
        \addlinespace
        Segurança adequada &
        Cap.~\ref{cap:seguranca} &
        BCrypt, JWT, roles, SecureStore, OWASP mobile \\
        \addlinespace
        Clareza na apresentação e domínio técnico &
        Cap.~\ref{cap:objetivos}, \ref{cap:conclusao} &
        Mapeamento de critérios; glossário implícito; decisões justificadas \\
        \addlinespace
        Estabilidade geral &
        Cap.~\ref{cap:testes} &
        Testes unitários backend; tratamento de erros; validações \\
        \bottomrule
    \end{tabular}
    \\[0.5em]
    \footnotesize Fonte: elaborado pelos autores com base no enunciado da disciplina e em \texttt{todo.md}.
\end{table}

\section{Mapeamento à metodologia e atividades da apresentação}

A Tabela~\ref{tab:atividades-apresentacao} indica como cada atividade exigida na metodologia será conduzida na banca.

\begin{table}[htbp]
    \caption{Atividades de apresentação e correspondência no Cardly}
    \label{tab:atividades-apresentacao}
    \centering
    \small
    \begin{tabular}{p{5.5cm} p{8.5cm}}
        \toprule
        \textbf{Atividade (professor)} & \textbf{Como será demonstrada} \\
        \midrule
        Proposta (problema e solução) & Slides iniciais + Cap.~1 deste relatório \\
        Fluxo completo (login, navegação, CRUD, API) & Demo no emulador/dispositivo: Login $\rightarrow$ Decks $\rightarrow$ Estudo $\rightarrow$ Dashboard \\
        API em funcionamento & Postman ou log ao vivo; endpoints \texttt{/api/auth}, \texttt{/api/decks}, \texttt{/api/cards} \\
        Arquitetura (Controller/Service/Repository/Hooks/Services) & Diagramas Cap.~\ref{cap:arquitetura}; navegação de pastas nos repositórios \\
        UI/UX e layout & FlipCard, paleta de cores, calendário no dashboard \\
        Build ou preparação para deploy & \texttt{eas build} ou APK pré-gerado; URL Railway ativa \\
        Questionamentos técnicos & Equipe dividida por domínio (ver roteiro de apresentação em material complementar) \\
        \bottomrule
    \end{tabular}
\end{table}

\section{Delimitações}

Não fazem parte do escopo mínimo desta entrega final de documentação: implementação completa do módulo de amizades (solicitações entre usuários), embora conste nos requisitos funcionais planejados; geração automática de slides \texttt{pptx}; e revisão por três agentes revisores exigida internamente pela equipe. Esses itens constam como pendências na conclusão deste relatório.

\chapter{Requisitos}
\label{cap:requisitos}

\section{Requisitos funcionais}

Os requisitos funcionais do MVP foram derivados do enunciado acadêmico, do documento \textit{Flashcards.pdf} e do arquivo \texttt{todo.md}. A Tabela~\ref{tab:rf-funcionais} consolida o estado de implementação conhecido na redação deste relatório.

\begin{longtable}{p{1.2cm} p{7.5cm} p{2.5cm} p{2.3cm}}
    \caption{Requisitos funcionais do Cardly} \label{tab:rf-funcionais} \\
    \toprule
    \textbf{ID} & \textbf{Descrição} & \textbf{Prioridade} & \textbf{Status} \\
    \midrule
    \endfirsthead
    \toprule
    \textbf{ID} & \textbf{Descrição} & \textbf{Prioridade} & \textbf{Status} \\
    \midrule
    \endhead
    RF01 & Criar conta de usuário & Alta & Implementado \\
    RF02 & Autenticar-se (login) & Alta & Implementado \\
    RF03 & Excluir própria conta & Alta & Implementado \\
    RF04 & CRUD de assuntos (\textit{decks}) & Alta & Implementado \\
    RF05 & Definir assunto público ou privado & Alta & Implementado \\
    RF06 & CRUD de cartas & Alta & Implementado \\
    RF07 & Responder carta (certo/errado) & Alta & Implementado \\
    RF08 & Pular carta na sessão de estudo & Alta & Implementado \\
    RF09 & Visualizar dashboard de progresso & Alta & Implementado \\
    RF10 & Revisar carta após intervalo & Alta & Implementado \\
    RF11 & Visualizar assuntos da comunidade & Média & Implementado \\
    RF12 & Clonar assunto público & Média & Implementado \\
    RF13 & Calendário de dias com estudo & Alta & Parcial \\
    RF14 & Solicitações de amizade & Baixa & Pendente \\
    RF15 & Admin: gerenciar usuários & Alta & Implementado \\
    RF16 & Admin: gerenciar assuntos e cartas & Alta & Implementado \\
    \bottomrule
    \multicolumn{4}{p{\linewidth}}{\footnotesize Fonte: elaborado pelos autores. Status sujeito a atualização conforme sprint final.}
\end{longtable}

\subsection{Regras de negócio: revisão espaçada}

O professor exige intervalo em dias entre tentativas conforme dificuldade da carta, calculada pelo sistema --- não pelo usuário \cite{anki_manual,anki_sm2}. O Cardly implementa níveis \texttt{NONE}, \texttt{EASY}, \texttt{MEDIUM} e \texttt{HARD}, com intervalos discretos de 1, 2, 3, 5, 7 e 9 dias \cite{medium_spaced_repetition}.

Regras principais:
\begin{itemize}
    \item Carta nova respondida corretamente torna-se \texttt{EASY} com revisão em 5 dias; incorreta torna-se \texttt{MEDIUM} com 3 dias.
    \item Em \texttt{EASY}, acertos consecutivos ampliam o intervalo até 9 dias; erros reduzem gradualmente o intervalo ou rebaixam para \texttt{MEDIUM}.
    \item Em \texttt{HARD}, erros reduzem intervalo para 1 dia; acertos promovem recuperação gradual para \texttt{MEDIUM}.
    \item Carta com \texttt{dueAt} futuro não pode ser respondida antes da data; após o vencimento, a resposta reagenda conforme desempenho.
\end{itemize}

A especificação completa encontra-se documentada internamente em \texttt{card-difficulty-level-rules.md} no repositório backend.

\subsection{Requisito de calendário}

O dashboard deve destacar dias em que o usuário respondeu cartas, com visual inspirado em aplicativos de controle de ciclo como o Clue \cite{clue_app}. A contagem baseia-se em registros de cartas atualizadas com intervalo agendado no dia. A interface visual do calendário permanece em refinamento (RF13 parcial).

\section{Requisitos técnicos}

\begin{longtable}{p{1.2cm} p{8cm} p{3.3cm}}
    \caption{Requisitos técnicos do MVP} \label{tab:rf-tecnicos} \\
    \toprule
    \textbf{ID} & \textbf{Descrição} & \textbf{Status} \\
    \midrule
    \endfirsthead
    RT01 & Autenticação JWT + Google SSO & Implementado \\
    RT02 & Estrutura de pastas backend padrão 1.0 & Implementado \\
    RT03 & Nomenclatura de arquivos padrão 1.1 & Implementado \\
    RT04 & Respostas paginadas (\texttt{PagedResponse}) & Implementado \\
    RT05 & Busca via \texttt{/search} + Specification & Implementado \\
    RT06 & Paleta de cores padrão 2.0 no front-end & Implementado \\
    RT07 & Modais de confirmação e toasts & Implementado \\
    RT08 & Estilização com NativeWind & Implementado \\
    RT09 & PostgreSQL + Flyway migrations & Implementado \\
    RT10 & Soft delete (\texttt{deletedAt}) em entidades & Implementado \\
    \bottomrule
\end{longtable}

\section{Requisitos não funcionais}

\begin{itemize}
    \item \textbf{Desempenho}: paginação padrão de 20 itens; consultas indexadas por usuário e deck.
    \item \textbf{Usabilidade}: feedback imediato via toast; confirmação antes de exclusões destrutivas \cite{mobile_first_nielsen}.
    \item \textbf{Manutenibilidade}: separação Controller/Service/Repository; DTOs imutáveis com \textit{records} Java \cite{java_records_spring}.
    \item \textbf{Portabilidade}: Expo permite builds Android/iOS; API independente de plataforma \cite{expo_eas_build}.
    \item \textbf{Segurança}: senhas com BCrypt; tokens com expiração; rotas admin restritas a \texttt{SUPERADMIN} \cite{bcrypt_spring,owasp_mobile}.
\end{itemize}

\section{Perfis de usuário}

Conforme requisitos acadêmicos institucionais, o sistema prevê perfis distintos:
\begin{description}
    \item[Estudante] Usuário padrão: gerencia próprios assuntos, estuda cartas, acessa comunidade e dashboard.
    \item[Superadministrador] Acessa \texttt{AdminScreen} exclusiva no app e endpoints \texttt{/api/admin/**} protegidos por role \cite{spring_security_jwt}.
\end{description}

\section{Rastreabilidade}

Cada requisito funcional mapeia-se a endpoints e telas específicas, detalhados nos Capítulos~\ref{cap:backend} e~\ref{cap:frontend}. Requisitos pendentes (RF13 completo, RF14) devem constar no plano de trabalhos futuros (Capítulo~\ref{cap:conclusao}).

\chapter{Arquitetura do Sistema}
\label{cap:arquitetura}

\section{Visão geral cliente-servidor}

O Cardly adota arquitetura cliente-servidor em três camadas lógicas \cite{tanenbaum2015,sommerville2019}:
\begin{enumerate}
    \item \textbf{Apresentação}: aplicativo React Native executado via Expo no dispositivo do estudante \cite{react_native_docs,expo_docs}.
    \item \textbf{Aplicação}: API REST Spring Boot expondo regras de negócio e autenticação \cite{spring_boot_docs,rest_fielding}.
    \item \textbf{Dados}: banco relacional PostgreSQL com schema versionado por Flyway \cite{postgresql_docs,flyway_docs}.
\end{enumerate}

A comunicação ocorre exclusivamente por HTTP/JSON sobre TLS em produção. O cliente não acessa o banco diretamente, garantindo controle centralizado de validações e agendamento de cartas \cite{clean_architecture_martin}.

\section{Diagrama de contexto}

\begin{figure}[htbp]
    \centering
    \fbox{\parbox{0.92\textwidth}{\centering
    \vspace{0.3cm}
    \textbf{[Figura 1 -- Diagrama de contexto]}\\
    \vspace{0.2cm}
    Dispositivo mobile (Expo) $\leftrightarrow$ API Spring Boot $\leftrightarrow$ PostgreSQL\\
    Integração opcional: Google OAuth 2.0\\
    \vspace{0.3cm}
    \textit{Inserir diagrama exportado do draw.io ou modelo institucional.}\\
    \vspace{0.3cm}
    }}
    \caption{Contexto do Cardly em ambiente de produção}
    \label{fig:contexto}
    \\[0.3em]
    \footnotesize Fonte: elaborado pelos autores.
\end{figure}

\section{Diagrama de camadas do back-end}

O backend segue o padrão de pastas 1.0 exigido pelo projeto:

\begin{itemize}
    \item \texttt{web} (Controller + DTO Request/Response): entrada HTTP, validação superficial, mapeamento de status.
    \item \texttt{service}: regras de negócio, transações, agendamento de cartas.
    \item \texttt{repository}: persistência JPA.
    \item \texttt{specification}: filtros dinâmicos para endpoints \texttt{/search}.
    \item \texttt{domain}: entidades JPA (\texttt{User}, \texttt{Deck}, \texttt{Card}).
    \item \texttt{config} e \texttt{security}: JWT, CORS, codificação de senha.
\end{itemize}

\begin{figure}[htbp]
    \centering
    \fbox{\parbox{0.92\textwidth}{\centering
    \vspace{0.3cm}
    \textbf{[Figura 2 -- Camadas Spring Boot]}\\
    Cliente HTTP $\rightarrow$ Controller $\rightarrow$ Service $\rightarrow$ Repository $\rightarrow$ BD\\
    Specification acoplada ao Repository em consultas paginadas\\
    \vspace{0.3cm}
    }}
    \caption{Fluxo em camadas do back-end Cardly}
    \label{fig:camadas-backend}
\end{figure}

\section{Diagrama de camadas do front-end}

O aplicativo mobile organiza responsabilidades assim \cite{react_hooks_docs}:

\begin{itemize}
    \item \texttt{screens}: telas de fluxo (Login, Decks, Estudo, Dashboard, Admin, Comunidade).
    \item \texttt{components}: UI reutilizável (\texttt{FlipCard}, \texttt{ConfirmModal}).
    \item \texttt{api}: funções de integração REST (\texttt{authApi}, \texttt{decksApi}, \texttt{dashboardApi}).
    \item \texttt{auth}: \texttt{AuthContext} com hooks de sessão e armazenamento seguro \cite{expo_secure_store}.
    \item \texttt{navigation}: pilhas autenticada e não autenticada \cite{react_navigation}.
    \item \texttt{theme}: tokens da paleta 2.0 consumidos pelo NativeWind \cite{nativewind_docs}.
\end{itemize}

\section{Fluxo de integração React Native + Spring Boot}

A integração completa --- critério central da disciplina --- segue o fluxo abaixo:

\subsection{Autenticação}

\begin{enumerate}
    \item Usuário informa credenciais ou token Google na tela de login.
    \item App envia \texttt{POST /api/auth/login} ou \texttt{POST /api/auth/google}.
    \item Backend valida credenciais (BCrypt) ou ID token Google \cite{google_oauth,bcrypt_spring}.
    \item Resposta inclui JWT assinado com claims de \texttt{userId} e \texttt{role} \cite{jwt_io}.
    \item App persiste token no \texttt{SecureStore} e injeta header \texttt{Authorization: Bearer} \cite{expo_secure_store,mdn_fetch}.
\end{enumerate}

\subsection{Operações CRUD}

\begin{enumerate}
    \item Tela dispara função em \texttt{*Api.ts} com token do contexto.
    \item Requisição atinge controller correspondente (\texttt{DeckController}, \texttt{CardController}).
    \item Service aplica autorização (usuário dono do recurso) e persistência.
    \item Resposta paginada (\texttt{PagedResponse}) ou DTO único retorna ao app.
    \item UI atualiza estado local e exibe toast de sucesso ou erro \cite{toast_message_rn}.
\end{enumerate}

\subsection{Sessão de estudo}

\begin{enumerate}
    \item \texttt{StudySessionScreen} busca cartas elegíveis (\texttt{dueOnly=true}).
    \item Usuário visualiza pergunta; toque aciona animação \texttt{FlipCard} \cite{w3schools_flip,reanimated_docs}.
    \item Resposta enviada via \texttt{POST /api/decks/\{id\}/cards/\{id\}/answer}.
    \item \texttt{CardService} recalcula dificuldade e \texttt{dueAt}; dashboard reflete \texttt{answeredToday}.
\end{enumerate}

\section{Contrato de API}

Endpoints principais (prefixo \texttt{/api}):

\begin{table}[htbp]
    \caption{Endpoints REST principais do Cardly}
    \label{tab:endpoints}
    \centering
    \small
    \begin{tabular}{lll}
        \toprule
        \textbf{Método} & \textbf{Rota} & \textbf{Descrição} \\
        \midrule
        POST & /auth/register, /auth/login, /auth/google & Autenticação \\
        GET & /me & Perfil autenticado \\
        POST & /decks/search & Busca paginada de assuntos \\
        POST & /decks/\{id\}/cards/search & Busca paginada de cartas \\
        POST & /decks/\{id\}/cards/\{id\}/answer & Resposta e agendamento \\
        GET & /dashboard & Indicadores agregados \\
        GET & /community/decks & Assuntos públicos \\
        POST & /community/decks/\{id\}/clone & Clonagem \\
        POST & /admin/users/search & Admin: usuários \\
        \bottomrule
    \end{tabular}
\end{table}

Documentação complementar pode ser gerada com springdoc-openapi \cite{openapi_springdoc}.

\section{Modelo de dados simplificado}

Entidades centrais:
\begin{itemize}
    \item \textbf{User}: email, senha hash, role (\texttt{USER} ou \texttt{SUPERADMIN}).
    \item \textbf{Deck}: título, visibilidade pública/privada, dono.
    \item \textbf{Card}: pergunta, resposta, \texttt{difficultyLevel}, \texttt{scheduledInterval}, \texttt{dueAt}, streaks.
\end{itemize}

Relacionamentos: User 1:N Deck; Deck 1:N Card. Soft delete via \texttt{deletedAt} em \texttt{BaseEntity}.

\section{Decisões arquiteturais}

\begin{itemize}
    \item \textbf{Stateless JWT}: escalabilidade horizontal do backend sem sessão servidor \cite{spring_security_jwt}.
    \item \textbf{Specification JPA}: atende requisito de \texttt{/search} sem proliferar métodos no repository \cite{jpa_specification}.
    \item \textbf{Expo managed workflow}: acelera entrega acadêmica; EAS Build para artefato instalável \cite{expo_eas_build}.
    \item \textbf{NativeWind}: unifica design tokens com utilitários Tailwind, reduzindo CSS inline \cite{nativewind_docs,tailwind_docs}.
\end{itemize}

\section{Preparação para produção na arquitetura}

Princípios twelve-factor \cite{twelve_factor} aplicados:
\begin{itemize}
    \item Configuração por variáveis de ambiente (\texttt{DATABASE\_URL}, \texttt{JWT\_SECRET}, \texttt{GOOGLE\_CLIENT\_ID}).
    \item Logs estruturados no backend (Spring Boot default).
    \item CORS restrito ao domínio do app \cite{helmet_cors_spring}.
    \item Migrations Flyway executadas no startup do container Railway.
    \item Front-end aponta \texttt{EXPO\_PUBLIC\_API\_URL} para instância publicada.
\end{itemize}

Detalhes operacionais no Capítulo~\ref{cap:deploy}.

\chapter{Implementação do Back-End}
\label{cap:backend}

\section{Stack tecnológica}

O back-end utiliza Java 17+, Spring Boot 3, Spring Data JPA, Spring Security, PostgreSQL e Flyway \cite{spring_boot_docs,postgresql_docs,flyway_docs}. O código-fonte está no repositório \texttt{cardly-backend} \cite{cardly_backend_repo}.

\section{Organização de pacotes (padrão 1.0 e 1.1)}

A estrutura segue o padrão exigido pelo professor:

\begin{verbatim}
com.cardly
├── config/          (SecurityConfig, JwtProperties)
├── domain/          (User, Deck, Card, ENUMs)
├── repository/      (UserRepository, DeckRepository, CardRepository)
├── service/         (AuthService, UserService, DeckService, CardService)
├── specification/   (UserSpecification, DeckSpecification, CardSpecification)
├── security/        (JwtAuthenticationFilter, UserPrincipal)
├── util/            (EmailNormalizer)
└── web/             (Controllers)
    └── dto/         (Request/Response records)
\end{verbatim}

Convenção de nomes: \texttt{UserController}, \texttt{UserService}, \texttt{UserRepository}, \texttt{UserSpecification}, \texttt{UserRoleENUM}, \texttt{RegisterRequest}, \texttt{UserResponse} \cite{youtube_spring_boot_api}.

\section{Camada de domínio}

\subsection{Entidade Card}

A entidade \texttt{Card} centraliza a lógica de revisão espaçada. Campos relevantes:
\begin{itemize}
    \item \texttt{difficultyLevel}: enum \texttt{NONE}, \texttt{EASY}, \texttt{MEDIUM}, \texttt{HARD}.
    \item \texttt{scheduledInterval}: enum \texttt{DAYS\_1} a \texttt{DAYS\_9}.
    \item \texttt{dueAt}: instante UTC da próxima revisão permitida.
    \item \texttt{rightStreak} e \texttt{wrongStreak}: contadores consecutivos.
\end{itemize}

\subsection{Soft delete}

Entidades estendem \texttt{BaseEntity} com \texttt{createdAt}, \texttt{updatedAt} e \texttt{deletedAt}. Specifications filtram \texttt{deletedAt IS NULL}, preservando integridade referencial e auditoria.

\section{Camada de serviço}

\subsection{CardService -- trecho crucial}

O método \texttt{answerCard} concentra as regras de negócio exigidas pelo docente. Ao receber \texttt{correct} ou \texttt{skipped}, o serviço:
\begin{enumerate}
    \item Valida propriedade do deck e existência da carta.
    \item Impede resposta antes de \texttt{dueAt} quando agendada.
    \item Aplica transições de dificuldade conforme especificação interna.
    \item Persiste novo \texttt{dueAt} via \texttt{Instant.now().plus(days, ChronoUnit.DAYS)}.
\end{enumerate}

\begin{lstlisting}[language=Java,caption={Trecho simplificado de agendamento em CardService},label={lst:card-schedule}]
private void applySchedule(Card card, DifficultyLevelENUM difficulty,
        ScheduledIntervalENUM interval) {
    card.setDifficultyLevel(difficulty);
    card.setScheduledInterval(interval);
    card.setDueAt(Instant.now().plus(intervalDays(interval), ChronoUnit.DAYS));
}
\end{lstlisting}

Referências teóricas: algoritmo SM-2 e adaptações práticas \cite{anki_sm2,medium_spaced_repetition}.

\subsection{DeckService e CommunityController}

Decks possuem flag \texttt{publicVisible}. Assuntos públicos aparecem em \texttt{/api/community/decks}; clonagem duplica deck e cartas para coleção do usuário, isolando alterações futuras do autor original.

\subsection{DashboardController}

Agrega métricas: total de assuntos, total de cartas, cartas vencidas (\texttt{dueCards}) e respostas no dia (\texttt{answeredToday}). Base para calendário e indicadores do app \cite{clue_app}.

\section{Camada web e DTOs}

Controllers expõem REST sem lógica de negócio. DTOs são \textit{records} Java imutáveis \cite{java_records_spring}:

\begin{itemize}
    \item \texttt{CardSearchRequest}: filtros \texttt{question}, \texttt{difficultyLevel}, \texttt{dueOnly}, paginação.
    \item \texttt{PagedResponse<T>}: \texttt{content}, \texttt{page}, \texttt{size}, \texttt{totalElements}, \texttt{totalPages}.
    \item \texttt{AnswerCardRequest}: \texttt{Boolean correct}, \texttt{Boolean skipped}.
\end{itemize}

Respostas de erro padronizadas em \texttt{ApiExceptionHandler}.

\section{Specification e endpoints /search}

Em substituição a \texttt{getAll}, cada recurso expõe \texttt{POST .../search} \cite{jpa_specification}:

\begin{lstlisting}[language=Java,caption={Composição de Specification para cartas},label={lst:spec}]
Specification<Card> spec = Specification.where(CardSpecification.active())
    .and(CardSpecification.byOwner(userId))
    .and(CardSpecification.byDeck(deckId))
    .and(CardSpecification.questionContains(request.question()))
    .and(CardSpecification.difficultyEquals(request.difficultyLevel()))
    .and(CardSpecification.dueNow(request.dueOnly()));
\end{lstlisting}

Payload vazio retorna conjunto paginado completo (filtrado por usuário autenticado), conforme requisito técnico RT05.

\section{Persistência e migrations}

Flyway scripts em \texttt{db/migration} criam tabelas \texttt{users}, \texttt{decks}, \texttt{cards} e índices por \texttt{user\_id} e \texttt{due\_at} \cite{flyway_docs}. Hibernate opera com \texttt{ddl-auto=validate} em produção.

\section{Autenticação no backend}

\texttt{AuthService} registra usuários com senha BCrypt, autentica login e valida ID token Google \cite{bcrypt_spring,google_oauth}. \texttt{JwtService} emite tokens HS256 com expiração configurável (\texttt{cardly.jwt.expiration-seconds}) \cite{jwt_io}.

\texttt{JwtAuthenticationFilter} intercepta requisições, valida assinatura e popula \texttt{SecurityContext} \cite{spring_security_jwt,stackoverflow_jwt_spring}.

\section{Administração}

\texttt{AdminController} restringe operações a \texttt{ROLE\_SUPERADMIN}. Endpoints espelham CRUD de usuários, decks e cartas com DTOs administrativos (\texttt{AdminCreateUserRequest}, \texttt{AdminCardSearchRequest}).

\section{Boas práticas adotadas}

\begin{itemize}
    \item \texttt{open-in-view=false} evita consultas lazy fora de transação.
    \item Serviços anotados com \texttt{@Transactional} delimitam consistência.
    \item Normalização de e-mail em \texttt{EmailNormalizer} reduz duplicatas.
    \item Paginação obrigatória evita respostas volumosas em mobile \cite{pressman2020}.
\end{itemize}

\section{Evidência para arguição}

Perguntas esperadas na banca e onde respondê-las:
\begin{itemize}
    \item \textit{Por que a dificuldade fica no backend?} --- Consistência e impossibilidade de fraude no agendamento.
    \item \textit{Por que Specification?} --- Requisito acadêmico e flexibilidade de filtros sem explosion de métodos.
    \item \textit{Como garantir paginação?} --- \texttt{Pageable} Spring Data em todos os \texttt{/search}.
\end{itemize}

\chapter{Implementação do Front-End}
\label{cap:frontend}

\section{Stack tecnológica}

O front-end utiliza React Native 0.81, Expo SDK 54, TypeScript 5.9, React Navigation 7, NativeWind 4 e React Native Reanimated \cite{react_native_docs,expo_docs,typescript_handbook,nativewind_docs,reanimated_docs}. Código em \texttt{cardly-rnative-app} \cite{cardly_frontend_repo}.

\section{Estrutura de pastas}

\begin{verbatim}
src/
├── api/           authApi, decksApi, dashboardApi, communityApi, adminApi
├── auth/          AuthContext.tsx
├── components/
│   ├── study/     FlipCard.tsx
│   └── ui/        ConfirmModal.tsx, toast.ts
├── navigation/    AuthStack, AppStack
├── screens/       Login, Register, Decks, Study, Dashboard, Admin...
└── theme/         colors.ts (paleta 2.0)
\end{verbatim}

Separação alinhada ao critério de avaliação sobre organização e responsabilidades \cite{clean_architecture_martin,react_hooks_docs}.

\section{Navegação e fluxos}

\subsection{Pilha de autenticação}

\texttt{AuthStack} contém \texttt{LoginScreen} e \texttt{RegisterScreen}. Após login bem-sucedido, \texttt{AuthContext} alterna para \texttt{AppStack}.

\subsection{Pilha principal}

\texttt{AppStack} agrupa:
\begin{itemize}
    \item \texttt{DashboardScreen}: hub com métricas e atalhos.
    \item \texttt{DecksScreen} / \texttt{DeckDetailScreen}: CRUD de assuntos e cartas.
    \item \texttt{StudySessionScreen}: sessão de estudo com flip e botões certo/errado/pular.
    \item \texttt{CommunityScreen}: assuntos públicos e clonagem.
    \item \texttt{AdminScreen}: visível apenas para \texttt{SUPERADMIN}.
\end{itemize}

React Navigation gerencia histórico e tipagem de parâmetros via \texttt{AppStackParamList} \cite{react_navigation}.

\section{Integração com API}

Cada módulo \texttt{*Api.ts} encapsula \texttt{fetch} \cite{mdn_fetch}:

\begin{enumerate}
    \item Monta URL base a partir de \texttt{EXPO\_PUBLIC\_API\_URL}.
    \item Injeta header \texttt{Authorization} quando autenticado.
    \item Serializa body JSON em POST/PUT/PATCH.
    \item Deserializa \texttt{PagedResponse} ou lança erro interpretado pela UI.
\end{enumerate}

Exemplo de fluxo no dashboard:

\begin{lstlisting}[language=Java,caption={Chamada tipada ao dashboard (TypeScript)},label={lst:dashboard-fetch}]
const data = await fetchDashboard(token);
setDashboard(data);
\end{lstlisting}

\section{Autenticação no cliente}

\texttt{AuthContext} provê \texttt{token}, \texttt{user}, \texttt{signIn}, \texttt{signOut} e \texttt{removeAccount}. Token persistido com \texttt{expo-secure-store} \cite{expo_secure_store}, atendendo recomendações OWASP para dados sensíveis em mobile \cite{owasp_mobile}.

Login Google envia ID token ao endpoint \texttt{/api/auth/google} \cite{google_oauth}.

\section{UI/UX e paleta de cores (padrão 2.0)}

Tokens definidos em \texttt{theme/colors.ts} e replicados no \texttt{tailwind.config.js} \cite{tailwind_docs}:

\begin{itemize}
    \item Primária: \texttt{\#2563EB} (ações principais).
    \item Fundo claro: \texttt{\#F8FAFC}; superfície: \texttt{\#FFFFFF}.
    \item Texto primário: \texttt{\#0F172A}; secundário: \texttt{\#64748B}.
    \item Feedback: success \texttt{\#22C55E}, warning \texttt{\#F59E0B}, error \texttt{\#EF4444}.
\end{itemize}

NativeWind aplica classes utilitárias (\texttt{className="bg-background-light"}) mantendo consistência visual \cite{nativewind_docs}.

\section{Componente FlipCard}

Requisito de UX: animação de virada ao toque (não hover), inspirada em flip cards CSS \cite{w3schools_flip}. Implementação com \texttt{Animated.spring} e rotação Y \cite{reanimated_docs}:

\begin{itemize}
    \item Estado \texttt{flipped} controlado pela tela de estudo.
    \item Frente exibe pergunta; verso exibe resposta.
    \item \texttt{backfaceVisibility: 'hidden'} evita artefatos visuais.
\end{itemize}

\section{Feedback ao usuário}

\subsection{Toasts}

\texttt{react-native-toast-message} exibe mensagens de sucesso e erro com moderação --- requisito de equilíbrio do professor \cite{toast_message_rn}. Helpers \texttt{showSuccess} e \texttt{showError} centralizam textos.

\subsection{Modais de confirmação}

Exclusões de conta, assuntos e cartas passam por \texttt{ConfirmModal}, reduzindo ações acidentais \cite{mobile_first_nielsen}.

\section{Tela de estudo}

\texttt{StudySessionScreen} orquestra:
\begin{enumerate}
    \item Carregamento de cartas vencidas (\texttt{dueOnly: true}).
    \item Exibição sequencial com \texttt{FlipCard}.
    \item Botões ``Errei'', ``Acertei'' e ``Pular'' mapeados a \texttt{AnswerCardRequest}.
    \item Toast e avanço automático para próxima carta.
\end{enumerate}

\section{Dashboard e calendário}

\texttt{DashboardScreen} exibe cards com \texttt{totalSubjects}, \texttt{totalCards}, \texttt{dueCards} e \texttt{answeredToday}. O calendário mensal com destaque de dias estudados está em implementação parcial (RF13); o design referencia aplicativos como Clue \cite{clue_app}.

\begin{figure}[htbp]
    \centering
    \fbox{\parbox{0.85\textwidth}{\centering
    \vspace{2cm}
    \textbf{[Figura 3 -- Print do Dashboard Cardly]}\\
    \vspace{2cm}
    }}
    \caption{Dashboard com indicadores de progresso (inserir captura de tela real)}
    \label{fig:dashboard-print}
\end{figure}

\section{Tela administrativa}

\texttt{AdminScreen} é renderizada condicionalmente quando \texttt{user.role === 'SUPERADMIN'}, atendendo requisito de front-end exclusivo para superadmin. Consome \texttt{adminApi} para buscas paginadas e formulários de criação/edição.

\section{Responsividade}

Layouts utilizam \texttt{flex-1}, padding consistente e tipografia escalável. Testes em emulador Android e dispositivo físico via Expo Go \cite{expo_docs,youtube_expo_typescript}.

\section{Hooks e estado}

Telas usam \texttt{useState}, \texttt{useEffect} e \texttt{useCallback} seguindo regras oficiais de Hooks \cite{react_hooks_docs}. Evita-se prop drilling excessivo delegando autenticação ao contexto global.

\section{Pendências de front-end}

\begin{itemize}
    \item Calendário visual completo no dashboard.
    \item Módulo de amizades (telas e integração API).
    \item Modo escuro completo com tokens \texttt{background.dark} (estrutura preparada na paleta).
\end{itemize}

\chapter{Segurança}
\label{cap:seguranca}

\section{Visão geral}

A segurança do Cardly combina autenticação forte, autorização baseada em papéis, comunicação protegida e boas práticas para aplicativos móveis \cite{owasp_mobile,spring_security_jwt,jwt_io}. Este capítulo atende ao critério de avaliação sobre implementação adequada de segurança.

\section{Autenticação}

\subsection{Credenciais locais}

Senhas nunca são armazenadas em texto claro. O backend utiliza \texttt{BCryptPasswordEncoder} com fator padrão Spring Security \cite{bcrypt_spring}. No registro, \texttt{AuthService} aplica hash antes de persistir \texttt{User.passwordHash}.

\subsection{JSON Web Token (JWT)}

Após login válido, o servidor emite JWT assinado com segredo configurável (\texttt{cardly.jwt.secret}). Claims incluem identificador do usuário e role. Expiração padrão: 3600 segundos (1 hora), configurável \cite{jwt_io,youtube_jwt_explained}.

O filtro \texttt{JwtAuthenticationFilter} valida token em cada requisição autenticada, mantendo API stateless \cite{spring_security_jwt,stackoverflow_jwt_spring}.

\subsection{Google Single Sign-On (SSO)}

Fluxo OAuth 2.0 \cite{google_oauth}:
\begin{enumerate}
    \item App obtém ID token via SDK Google/Expo (configuração \texttt{cardly.google.client-id}).
    \item Backend valida audiência e emissor do token.
    \item Usuário existente autentica; novo e-mail pode gerar registro automático conforme regra de negócio.
\end{enumerate}

\section{Autorização}

\texttt{SecurityConfig} define regras:

\begin{itemize}
    \item \texttt{/api/auth/**}: público.
    \item \texttt{/api/admin/**}: exige \texttt{ROLE\_SUPERADMIN}.
    \item Demais rotas: autenticadas.
\end{itemize}

Serviços validam ownership: usuário comum só manipula decks/cartas próprios. Admin bypassa via endpoints dedicados auditáveis.

\section{Proteção no cliente mobile}

\begin{itemize}
    \item Token JWT armazenado em \texttt{SecureStore}, não em AsyncStorage plain \cite{expo_secure_store}.
    \item Logout remove token local e invalida sessão no contexto.
    \item Telas admin ocultas por role, não apenas por rota deep link (defesa em profundidade com backend).
\end{itemize}

\section{Comunicação e CORS}

Em produção, API exposta via HTTPS (Railway/AWS). CORS configurado para origens confiáveis do app e ferramentas de teste \cite{helmet_cors_spring}. CSRF desabilitado por API stateless JWT --- decisão documentada e comum em SPAs/mobile \cite{spring_security_jwt}.

\section{Validação e tratamento de erros}

\begin{itemize}
    \item DTOs validados com Bean Validation (\texttt{@NotBlank}, \texttt{@Email}).
    \item \texttt{ApiExceptionHandler} retorna HTTP semântico (400, 401, 403, 404) sem vazar stack trace.
    \item Mensagens genéricas em falha de login (``credenciais inválidas'') mitigam enumeração de usuários.
\end{itemize}

\section{Segredos e configuração}

Segredos não são versionados no Git. Em produção:
\begin{itemize}
    \item \texttt{JWT\_SECRET} com no mínimo 256 bits para HS256.
    \item Credenciais PostgreSQL via variáveis do provedor \cite{twelve_factor,railway_docs}.
    \item \texttt{GOOGLE\_CLIENT\_ID} distinto por ambiente (dev/prod).
\end{itemize}

\section{Soft delete e integridade}

Exclusões lógicas preservam histórico e impedem reutilização indevida de IDs. Admin pode restaurar ou purgar conforme política futura.

\section{Checklist OWASP Mobile (parcial)}

\begin{table}[htbp]
    \caption{Checklist de segurança mobile aplicado ao Cardly}
    \label{tab:owasp}
    \centering
    \small
    \begin{tabular}{p{5cm} p{8cm}}
        \toprule
        \textbf{Controle} & \textbf{Implementação} \\
        \midrule
        Armazenamento seguro de credenciais & SecureStore para JWT \\
        Transporte criptografado & HTTPS em produção \\
        Autenticação server-side & JWT validado no backend \\
        Controle de acesso & Roles + ownership nos services \\
        Validação de entrada & Bean Validation + tipos TS \\
        \bottomrule
    \end{tabular}
\end{table}

Referência: OWASP MASVS \cite{owasp_mobile}.

\section{Riscos residuais}

\begin{itemize}
    \item Segredo JWT de desenvolvimento no \texttt{application.properties} local --- deve ser substituído em produção.
    \item Rate limiting não implementado (aceitável para demo acadêmica; recomendado em produção real).
    \item Refresh token ausente; usuário reautentica após expiração de 1h.
\end{itemize}

\chapter{Testes e Validação}
\label{cap:testes}

\section{Estratégia de testes}

A validação do Cardly combina testes automatizados no backend, testes manuais integrados mobile-API e evidências visuais exigidas pelo manual acadêmico UNILESTE \cite{junit5_spring,jest_testing_library,postman_docs}. A estratégia prioriza regras de negócio críticas (autenticação, agendamento de cartas, autorização admin).

\section{Testes automatizados -- backend}

Suite JUnit 5 + Spring Boot Test executada via Maven (\texttt{mvn test}) \cite{junit5_spring}. Classes principais:

\begin{table}[htbp]
    \caption{Testes unitários e de integração do backend}
    \label{tab:testes-backend}
    \centering
    \small
    \begin{tabular}{ll}
        \toprule
        \textbf{Classe} & \textbf{Escopo} \\
        \midrule
        AuthServiceTest & Registro, login, hash de senha \\
        UserServiceTest & CRUD usuário, soft delete \\
        DeckServiceTest & Assuntos, visibilidade pública \\
        CardServiceTest & Respostas, intervalos, dueAt \\
        EmailNormalizerTest & Normalização de e-mail \\
        CardlyBackendApplicationTests & Context load Spring \\
        \bottomrule
    \end{tabular}
\end{table}

Ambiente de teste utiliza H2 ou PostgreSQL em container via \texttt{application-test.properties} \cite{docker_postgres}.

\subsection{Exemplo: validação de agendamento}

Testes em \texttt{CardServiceTest} verificam:
\begin{itemize}
    \item Carta nova correta $\rightarrow$ EASY, 5 dias.
    \item Carta nova incorreta $\rightarrow$ MEDIUM, 3 dias.
    \item Bloqueio de resposta antes de \texttt{dueAt}.
    \item Skip reagenda conforme dificuldade atual.
\end{itemize}

\section{Testes manuais -- API}

Postman (ou Insomnia) valida contratos REST \cite{postman_docs}:

\begin{enumerate}
    \item Registrar usuário e obter JWT.
    \item Criar deck e cartas.
    \item Responder carta e consultar novo \texttt{dueAt}.
    \item Tentar acessar \texttt{/api/admin} com role USER (esperado 403).
    \item Buscar decks com payload vazio em \texttt{/search} (paginação).
\end{enumerate}

pgAdmin inspeciona persistência no PostgreSQL \cite{pgadmin_docs}.

\section{Testes manuais -- aplicativo mobile}

Roteiro de validação integrada:

\begin{longtable}{p{1cm} p{5cm} p{7.5cm}}
    \caption{Roteiro de testes manuais mobile} \label{tab:testes-mobile} \\
    \toprule
    \textbf{ID} & \textbf{Cenário} & \textbf{Resultado esperado} \\
    \midrule
    \endfirsthead
    T01 & Registro e login & Token salvo; navegação ao dashboard \\
    T02 & Criar assunto privado & Lista atualizada; toast sucesso \\
    T03 & Adicionar cartas & Cartas visíveis no deck \\
    T04 & Sessão de estudo + flip & Animação; resposta registrada \\
    T05 & Carta agendada futura & Indisponível até dueAt \\
    T06 & Dashboard & Contadores coerentes com API \\
    T07 & Comunidade + clone & Deck copiado na coleção \\
    T08 & Admin (superadmin) & CRUD global funcional \\
    T09 & Excluir conta & Modal confirma; logout automático \\
    \bottomrule
\end{longtable}

\section{Evidências visuais (prints)}

O relatório final deve incluir capturas de tela de diversos cenários reais, conforme requisito institucional. Placeholders abaixo indicam prints pendentes de inserção:

\begin{figure}[htbp]
    \centering
    \fbox{\parbox{0.4\textwidth}{\centering\vspace{2.5cm}Login\vspace{2.5cm}}
    \fbox{\parbox{0.4\textwidth}{\centering\vspace{2.5cm}Estudo + Flip\vspace{2.5cm}}
    \caption{Evidências visuais -- login e sessão de estudo (inserir PNGs reais)}
    \label{fig:prints-estudo}
\end{figure}

\begin{figure}[htbp]
    \centering
    \fbox{\parbox{0.4\textwidth}{\centering\vspace{2.5cm}CRUD Decks\vspace{2.5cm}}
    \fbox{\parbox{0.4\textwidth}{\centering\vspace{2.5cm}Admin Panel\vspace{2.5cm}}
    \caption{Evidências visuais -- CRUD e administração}
    \label{fig:prints-crud}
\end{figure}

Recomenda-se pasta \texttt{docs/report/figures/} com arquivos \texttt{login.png}, \texttt{study.png}, etc.

\section{Testes de front-end (planejado)}

React Native Testing Library \cite{jest_testing_library} pode cobrir:
\begin{itemize}
    \item Renderização de \texttt{FlipCard} nos estados frente/verso.
    \item Exibição condicional de \texttt{AdminScreen}.
    \item Helpers de toast.
\end{itemize}

Cobertura automatizada mobile permanece como melhoria futura; validação atual é predominantemente manual integrada.

\section{CI/CD (opcional)}

GitHub Actions pode executar \texttt{mvn test} a cada push \cite{github_actions}. Pipeline não bloqueia entrega acadêmica, mas reforça estabilidade citada nos critérios de avaliação.

\section{Rastreabilidade requisito-teste}

\begin{table}[htbp]
    \caption{Rastreabilidade entre requisitos e testes}
    \label{tab:rastreabilidade}
    \centering
    \footnotesize
    \begin{tabular}{lll}
        \toprule
        \textbf{RF/RT} & \textbf{Teste} & \textbf{Tipo} \\
        \midrule
        RF01--RF03 & AuthServiceTest, T01, T09 & Auto + manual \\
        RF04--RF06 & DeckServiceTest, T02--T03 & Auto + manual \\
        RF07--RF10 & CardServiceTest, T04--T05 & Auto + manual \\
        RF09 & DashboardController, T06 & Auto + manual \\
        RF15--RF16 & Admin endpoints, T08 & Manual \\
        RT01 & AuthServiceTest, T01 & Auto + manual \\
        RT05 & CardServiceTest search, Postman & Auto + manual \\
        \bottomrule
    \end{tabular}
\end{table}

\section{Critérios de aceite para demonstração}

Para a apresentação prática, considera-se aceito:
\begin{itemize}
    \item Login $\rightarrow$ estudo $\rightarrow$ dashboard sem erro fatal.
    \item API respondendo em produção ou emulador apontando para Railway.
    \item Pelo menos um teste automatizado de carta demonstrado na arguição.
    \item Prints inseridos na versão final deste PDF.
\end{itemize}

\chapter{Deploy e Produção}
\label{cap:deploy}

\section{Objetivo}

O professor exige preparação para ambiente de produção, mesmo com operação temporária para apresentação \cite{twelve_factor,todo.md}. Este capítulo descreve arquitetura de deploy, variáveis de ambiente e procedimentos de build.

\section{Topologia de produção proposta}

\begin{figure}[htbp]
    \centering
    \fbox{\parbox{0.92\textwidth}{\centering
    \vspace{0.4cm}
    \textbf{[Figura 4 -- Topologia de deploy]}\\[0.3cm]
    Dispositivo mobile $\rightarrow$ HTTPS $\rightarrow$ Railway (Spring Boot + PostgreSQL)\\
    Expo EAS Build $\rightarrow$ APK/AAB instalável\\
    Opcional: Vercel para landing ou painel web auxiliar\\
    \vspace{0.4cm}
    }}
    \caption{Infraestrutura de produção do Cardly}
    \label{fig:deploy-topologia}
\end{figure}

Componentes \cite{railway_docs,vercel_docs,aws_free_tier,expo_eas_build}:
\begin{itemize}
    \item \textbf{Backend}: container Java executando JAR Spring Boot.
    \item \textbf{Banco}: PostgreSQL gerenciado (Railway plugin ou AWS RDS free tier).
    \item \textbf{Mobile}: binário Android via EAS Build; Expo Go apenas em desenvolvimento.
    \item \textbf{Front web opcional}: Vercel para documentação ou página de download.
\end{itemize}

\section{Configuração por ambiente}

Princípio twelve-factor: configuração via ambiente, não via código \cite{twelve_factor}.

\subsection{Backend (Railway)}

Variáveis obrigatórias:

\begin{table}[htbp]
    \caption{Variáveis de ambiente do backend em produção}
    \label{tab:env-backend}
    \centering
    \small
    \begin{tabular}{ll}
        \toprule
        \textbf{Variável} & \textbf{Descrição} \\
        \midrule
        SPRING\_DATASOURCE\_URL & JDBC PostgreSQL \\
        SPRING\_DATASOURCE\_USERNAME & Usuário BD \\
        SPRING\_DATASOURCE\_PASSWORD & Senha BD \\
        CARDLY\_JWT\_SECRET & Segredo HS256 (256+ bits) \\
        CARDLY\_JWT\_EXPIRATION\_SECONDS & TTL token (ex.: 3600) \\
        CARDLY\_GOOGLE\_CLIENT\_ID & Client ID OAuth Google \\
        SERVER\_PORT & Porta exposta (Railway injeta) \\
        \bottomrule
    \end{tabular}
\end{table}

Flyway executa migrations no startup \cite{flyway_docs}. Health check em \texttt{/actuator/health} (se habilitado).

\subsection{Frontend mobile (Expo)}

\begin{table}[htbp]
    \caption{Variáveis públicas do app Expo}
    \label{tab:env-frontend}
    \centering
    \small
    \begin{tabular}{ll}
        \toprule
        \textbf{Variável} & \textbf{Descrição} \\
        \midrule
        EXPO\_PUBLIC\_API\_URL & URL base HTTPS da API Railway \\
        EXPO\_PUBLIC\_GOOGLE\_CLIENT\_ID & Client ID para login Google \\
        \bottomrule
    \end{tabular}
\end{table}

Rebuild obrigatório após alterar variáveis \texttt{EXPO\_PUBLIC\_*} \cite{expo_docs}.

\section{Pipeline de deploy do backend}

Procedimento resumido \cite{railway_docs}:
\begin{enumerate}
    \item Conectar repositório GitHub \texttt{cardly-backend} ao Railway.
    \item Adicionar serviço PostgreSQL e vincular \texttt{DATABASE\_URL}.
    \item Configurar variáveis da Tabela~\ref{tab:env-backend}.
    \item Deploy automático a cada push na branch principal.
    \item Validar logs: Flyway migrations OK, Tomcat started.
\end{enumerate}

Alternativa AWS Free Tier \cite{aws_free_tier}: EC2 t2.micro + RDS PostgreSQL, com maior overhead operacional.

\section{Build mobile para demonstração}

\begin{enumerate}
    \item Instalar EAS CLI: \texttt{npm install -g eas-cli}.
    \item Configurar \texttt{eas.json} com profile \texttt{preview} (APK).
    \item Executar \texttt{eas build -p android --profile preview} \cite{expo_eas_build}.
    \item Distribuir APK via link Expo ou USB na apresentação.
\end{enumerate}

Para iOS, necessário Apple Developer (opcional no escopo Android-first da equipe).

\section{CORS e rede}

Backend deve permitir origem do app. Em APIs mobile puras, CORS afeta menos chamadas nativas, mas permanece relevante para testes web e ferramentas \cite{helmet_cors_spring}. Recomenda-se lista explícita de origens, não wildcard em produção.

\section{Monitoramento e logs}

Railway exibe logs stdout/stderr. Métricas básicas: latência, taxa de 5xx, uso de memória JVM. Para demo acadêmica, monitoramento avançado (Prometheus/Grafana) é opcional.

\section{Plano de rollback}

\begin{itemize}
    \item Railway permite redeploy de release anterior.
    \item Migrations Flyway devem ser backward-compatible ou acompanhadas de scripts de reversão.
    \item APK anterior mantido localmente como fallback na apresentação.
\end{itemize}

\section{Checklist pré-apresentação}

\begin{enumerate}
    \item API pública respondendo \texttt{GET /api/me} com token válido.
    \item Banco populado com dados de demonstração (usuário admin + decks exemplo).
    \item APK compilado com \texttt{EXPO\_PUBLIC\_API\_URL} de produção.
    \item Segredo JWT de produção diferente do desenvolvimento.
    \item Google OAuth configurado com SHA-1 do keystore de release.
    \item Teste completo T01--T09 (Capítulo~\ref{cap:testes}) em rede móvel (4G/Wi-Fi).
\end{enumerate}

\section{Descomissionamento}

Conforme escopo acadêmico, após apresentação a equipe pode:
\begin{itemize}
    \item Desativar serviço Railway e banco.
    \item Revogar client IDs OAuth de teste.
    \item Manter repositórios públicos como portfólio.
\end{itemize}

\section{Production readiness -- síntese}

\begin{table}[htbp]
    \caption{Matriz de prontidão para produção}
    \label{tab:production-readiness}
    \centering
    \footnotesize
    \begin{tabular}{p{3.5cm} p{2cm} p{7cm}}
        \toprule
        \textbf{Aspecto} & \textbf{Status} & \textbf{Observação} \\
        \midrule
        Config externalizada & OK & Env vars Railway/Expo \\
        Migrations versionadas & OK & Flyway \\
        Autenticação segura & Parcial & Falta refresh token \\
        HTTPS & OK & Railway default \\
        Build reprodutível & OK & Maven + EAS \\
        Observabilidade & Básico & Logs Railway \\
        Rate limiting & Pendente & Aceitável para demo \\
        Backup BD & Manual & Snapshot Railway pré-demo \\
        \bottomrule
    \end{tabular}
\end{table}

Referências: The Twelve-Factor App \cite{twelve_factor}; documentação Railway e Expo \cite{railway_docs,expo_eas_build}.

\chapter{Conclusão}
\label{cap:conclusao}

\section{Síntese do trabalho}

O Cardly consolidou um ecossistema de estudo mobile integrado a API REST, atendendo ao objetivo central da disciplina: demonstrar domínio da integração React Native (Expo/TypeScript) com Spring Boot (Java) em arquitetura cliente-servidor \cite{react_native_docs,spring_boot_docs,tanenbaum2015}.

Principais entregas técnicas:
\begin{itemize}
    \item Autenticação JWT com suporte a Google SSO e perfil superadministrador \cite{jwt_io,google_oauth}.
    \item CRUD completo de assuntos e cartas com busca paginada via Specification \cite{jpa_specification}.
    \item Regras de revisão espaçada implementadas no \texttt{CardService}, com intervalos de 1 a 9 dias \cite{anki_manual,medium_spaced_repetition}.
    \item Front-end com NativeWind, flip animado, toasts e modais \cite{nativewind_docs,w3schools_flip}.
    \item Preparação para deploy em Railway e build mobile via EAS \cite{railway_docs,expo_eas_build}.
\end{itemize}

\section{Atendimento aos critérios de avaliação}

Conforme mapeamento do Capítulo~\ref{cap:objetivos} (Tabela~\ref{tab:criterios-avaliacao}), o projeto cobre:
\begin{itemize}
    \item Integração front-back: evidenciada nos capítulos de arquitetura, backend e frontend.
    \item Segurança: JWT, BCrypt, roles e SecureStore (Capítulo~\ref{cap:seguranca}).
    \item Organização de código: padrões 1.0, 1.1 e 2.0 documentados.
    \item UX: sessão de estudo fluida, paleta consistente, feedback ao usuário.
    \item Produção: variáveis de ambiente, migrations e checklist de deploy.
\end{itemize}

\section{Limitações identificadas}

\begin{enumerate}
    \item Calendário visual no dashboard ainda parcial (RF13).
    \item Módulo de amizades não implementado (RF14).
    \item Ausência de refresh token JWT; sessão expira em uma hora.
    \item Cobertura automatizada limitada no front-end mobile.
    \item Rate limiting e monitoramento avançado não implementados.
\end{enumerate}

Limitações são aceitáveis no recorte de MVP acadêmico, mas devem ser explicitadas na arguição.

\section{Trabalhos futuros}

\begin{itemize}
    \item Concluir calendário estilo Clue com heatmap mensal \cite{clue_app}.
    \item Implementar solicitações de amizade e compartilhamento social de decks.
    \item Adicionar notificações push (Expo Notifications) para cartas vencidas.
    \item Expandir testes E2E com Detox ou Maestro.
    \item Publicar OpenAPI via springdoc para documentação interativa \cite{openapi_springdoc}.
    \item Modo escuro completo utilizando tokens \texttt{background.dark}.
    \item Sincronização offline com fila de respostas (desafio avançado mobile).
\end{itemize}

\section{Aprendizados da equipe}

Estudantes do sétimo semestre reportaram curva de aprendizado significativa em:
\begin{itemize}
    \item Configuração de JWT end-to-end \cite{youtube_jwt_explained,stackoverflow_jwt_spring}.
    \item Animações nativas com Reanimated \cite{reanimated_docs,youtube_expo_typescript}.
    \item Specification JPA para filtros dinâmicos \cite{jpa_specification}.
    \item Disciplina de variáveis de ambiente e deploy \cite{twelve_factor,railway_docs}.
\end{itemize}

Vídeos e artigos citados ao longo deste relatório \cite{youtube_spring_boot_api,medium_spaced_repetition} foram recursos pedagógicos essenciais.

\section{Considerações finais}

O Cardly demonstra que flashcards com revisão espaçada constituem domínio adequado para aplicação mobile com regras de negócio não triviais, indo além de CRUD superficial exigido em projetos introdutórios. A documentação formal em ABNT \cite{abnt_nbr14724,abnt_nbr6023}, aliada à demonstração funcional integrada, prepara a equipe para apresentação prática e arguição técnica previstas na metodologia da disciplina.

A versão final deste relatório depende da inserção de figuras reais (prints, diagramas), revisão ortográfica por três revisores (dois com modelo Opus conforme planejamento interno), dados do orientador na capa e compilação PDF via Overleaf ou LaTeX local com fonte Times New Roman conforme manual UNILESTE.

\section{Referências aos repositórios}

Código-fonte e histórico de desenvolvimento:
\begin{itemize}
    \item Backend: \url{https://github.com/hugoFreit4s/cardly-backend} \cite{cardly_backend_repo}
    \item Frontend: \url{https://github.com/hugoFreit4s/cardly-rnative-app} \cite{cardly_frontend_repo}
\end{itemize}


\cleardoublepage
\bibliography{references}

\end{document}
